% !TeX spellcheck = en_US
\documentclass[french]{yLectureNote}

\title{Outils Mathématiques 2}
\subtitle{Physique}
\author{Paulhenry Saux}
\date{\today}
\yLanguage{Français}

\professor{F.Pettinari}
\usepackage{graphicx}%----pour mettre des images
\usepackage[utf8]{inputenc}%---encodage
\usepackage{geometry}%---pour modifier les tailles et mettre a4paper
%\usepackage{awesomebox}%---pour les boites d'exercices, de pbq et de croquis ---d\'esactiv\'e pour les TP de PC
\usepackage{tikz}%---pour deiffner + d\'ependance de chemfig
% \usepackage{tabularx}%---pour dimensionner automatiquement les tableaux avec variable X
\usepackage{awesomebox}%---Pour les boites info, danger et autres
\usepackage{menukeys}%---Pour deiffner les touches de Calculatrice
\usepackage{fancyhdr}%---pour les en-t\^ete personnalis\'ees
\usepackage{blindtext}%---pour les liens
\usepackage{hyperref}%---pour les liens (\`a mettre en dernier)
\usepackage{caption}%---pour la francisation de la l\'egende table vers Tableau
\usepackage{pifont}
\usepackage{array}%---pour les tableaux
\usepackage{lipsum}
\usepackage{yFlatTable}
\usepackage{multicol}
\newcommand{\Lim}[1]{\lim\limits_{\substack{#1}}\:}
\renewcommand{\vec}{\overrightarrow}
\newcommand{\N}[0]{\mathbb{N}}
\newcommand{\dd}{\mathrm{d}}
\newcommand{\norm}[1]{||\vec{#1}||}
\newcommand{\fo}{\psi(\vec{r},t)}
\newcommand{\foe}{\psi(\vec{r},t)\*}
\newcommand{\HH}{\hat{H}}
\newcommand{\hb}{\hbar}
\newcommand{\lap}{\nabla^2}
\newcommand{\lapcc}{\frac{\partial^2 }{\partial x^2}+\frac{\partial^2 }{\partial y^2}+\frac{\partial^2 }{\partial z^2}}
\newcommand{\mpsi}{\(\psi\)}
\begin{document}
\setcounter{chapter}{2}
\chapter{Calcul vectoriel}
\section{Flux d'un champ de vecteur}
\begin{definition}[Flux d'un champ de vecteur]
\[\iint_S \vec{F}(\vec{r})\cdot \dd \vec{S} = \iint_S \dd S \vec{F}(\vec{r})\cdot \vec{n}\]
\end{definition}
\section{Gradient}
\begin{definition}[Gradient]
\[\nabla f = \frac{\partial }{\partial x}\vec{e_x}+\frac{\partial }{\partial y}\vec{e_y}+\frac{\partial }{\partial z}\vec{e_z}\]
\end{definition}
\begin{theorem}[Plan tangent à une surface]
L'équation est, avec un point \(M_0\) appartenant  à la surface : \[\nabla f(x_0,y_0,z_0)\cdot (x-x_0, y-y_0,z-z_0) = 0\]
\end{theorem}
\begin{theorem}[Plan tangent à une surface de la forme \(z = h(x,y)\)]
\[z=z_0+(x-x_0)\frac{\partial}{\partial x}(x_0,y_0)+(y-y_0)\frac{\partial h}{\partial y}(x_0,y_0)\]
\end{theorem}

\begin{theorem}[Condition de conservation]
En écrivant \(\vec{F} = P\vec{e_x}+ Q\vec{e_y}\) :
Pour un champ 2D : \(\frac{\partial P}{\partial y} = \frac{\partial Q}{\partial x}\)

Pour un champ 3D : \(\frac{\partial P}{\partial y} = \frac{\partial Q}{\partial Q}\),  \(\frac{\partial Q}{\partial z} = \frac{\partial R}{\partial y}\), \(\frac{\partial R}{\partial x} = \frac{\partial P}{\partial z}\)
\end{theorem}
\section{Autres}
\begin{definition}[Rotationnel]
\[rot(\vec{F}) = (\frac{\partial R}{\partial y}-\frac{\partial Q}{\partial z}),  (\frac{\partial P}{\partial z}-\frac{\partial R}{\partial x}),  (\frac{\partial Q}{\partial x}-\frac{\partial P}{\partial y})\]
\end{definition}
\begin{theorem}[Formule de Green]
 \[\int_C \vec{F}\cdot \dd \vec{r} = \iint_S \dd x\dd y (\frac{\partial Q}{\partial x}-\frac{\partial P}{\partial y})\]
\end{theorem}
\begin{definition}[Divergent]
\[dif(\vec{F})(x,y,z) = \frac{\partial P}{\partial x}(x,y,z) + \frac{\partial Q}{\partial y}(x,y,z) + \frac{\partial R}{\partial z}(x,y,z)\]
\end{definition}
\section{Méthode}
\subsection{Calculer le flux d'un champ de vecteur}
On cherche à calculer le flux de $\vec{F}(x,y,z)=\frac{1}{1+(z/l)^{2}}(x\vec{e}_{x}+y\vec{e}_{y}+z\vec{e}_{z})$ à travers la surface latérale $\mathcal{S}$ du cylindre infini d'axe $(O, \vec{e}_z)$ et de rayon $R$.
\subsubsection{Élément de surface orienté ($d\vec{S}$)}
Sur le cylindre de rayon $R$, $\dd\vec{S} = \vec{e}_\rho \, \dd S = \vec{e}_\rho \, R \, \dd\varphi \, dz$.
\subsubsection{Produit scalaire}
    $$\vec{F} \cdot \dd\vec{S} = \frac{1}{1+(z/l)^{2}} (x\vec{e}_{x}+y\vec{e}_{y}+z\vec{e}_{z}) \cdot \vec{e}_\rho \, R \, \dd\varphi \, \dd z$$

    En coordonnées cylindriques, $x\vec{e}_{x}+y\vec{e}_{y} = \rho \vec{e}_\rho$.

    Sur le cylindre, $\rho=R$.

    On en déduit
    $$(x\vec{e}_{x}+y\vec{e}_{y}+z\vec{e}_{z}) = R\vec{e}_{\rho} + z\vec{e}_{z}$$

    Le produit scalaire donne
    $$(R\vec{e}_{\rho} + z\vec{e}_{z}) \cdot \vec{e}_\rho = R(\vec{e}_{\rho}\cdot\vec{e}_{\rho}) + z(\vec{e}_{z}\cdot\vec{e}_{\rho}) = R$$
    $$\vec{F} \cdot \dd\vec{S} = \frac{R}{1+(z/l)^{2}} \, R \, \dd\varphi \, \dd z = \frac{R^2}{1+(z/l)^{2}} \, \dd\varphi \, \dd z$$
\subsubsection{Calcul du flux}
L'intégration se fait sur $z \in ]-\infty, +\infty[$ et $\varphi \in [0, 2\pi]$.
    $$\Phi = R^2 \int_{0}^{2\pi} d\varphi \cdot \int_{-\infty}^{+\infty} \frac{1}{1+(z/l)^{2}} \, dz$$
    $$\Phi = 2\pi R^2 \cdot (l [\arctan(z/l)]_{-\infty}^{+\infty}) = 2\pi R^2 \cdot (l (\frac{\pi}{2} - (-\frac{\pi}{2}))) =  2\pi^2 l R^2$$
\subsection{Calculer un plan tangent}
\subsubsection{Plan Tangent à $z=f(x,y) = x^2 + xy + y^2$ au point $(1, 2, 7)$}

L'équation du plan tangent à la surface $z=f(x,y)$ au point $(x_0, y_0, z_0)$ est donnée par :
$$
z - z_0 = \frac{\partial f}{\partial x}(x_0, y_0)(x - x_0) + \frac{\partial f}{\partial y}(x_0, y_0)(y - y_0)
$$

A. Calcul des dérivées partielles de $f(x,y) = x^{2}+x y+y^{2}$.
    $$
    \frac{\partial f}{\partial x} = 2x + y,  \frac{\partial f}{\partial y} = x + 2y
    $$

B. Évaluation au point $M_0 = (1, 2, 7)$

On a

    $$
    \frac{\partial f}{\partial x}(1, 2) = 2(1) + 2 = 4, \frac{\partial f}{\partial y}(1, 2) = 1 + 2(2) = 5
    $$
C. Équation du plan tangent

En substituant les valeurs dans l'équation du plan, on obtient :
\begin{flalign*}
z - 7 &= 4(x - 1) + 5(y - 2)\\
z - 7 &= 4x - 4 + 5y - 10\\
4x + 5y - z &= 4 + 10 - 7\\
4x + 5y - z &= 7\\
\end{flalign*}
\subsubsection{Plan Tangent à la surface implicite $f(x,y,z)=0$ au point $M_0=(1, 1, 1)$}

Ici, l'équation du plan tangent est donnée par :
$$
\vec{\nabla}f(M_0) \cdot (\vec{r} - \vec{r}_0) = 0
$$
soit
$$
\frac{\partial f}{\partial x}(M_0)(x - x_0) + \frac{\partial f}{\partial y}(M_0)(y - y_0) + \frac{\partial f}{\partial z}(M_0)(z - z_0) = 0
$$

La fonction est $f(x,y,z)=2x^{3}-3x^{2}+y^{2}z+yz^{2}+\frac{1}{3}z^{3}-\frac{4}{3}$.

A. Calcul des dérivées partielles :
    $$
    \frac{\partial f}{\partial x} = 6x^2 - 6x, \frac{\partial f}{\partial y} = 2yz + z^2,  \frac{\partial f}{\partial z} = y^2 + 2yz + z^2
    $$

B. Évaluation au point $M_0 = (1, 1, 1)$ :
    $$
    \frac{\partial f}{\partial x}(1, 1, 1) =  0, \frac{\partial f}{\partial y}(1, 1, 1) = 3, \frac{\partial f}{\partial z}(1, 1, 1) = 4
    $$
Le vecteur normal est donc $\vec{\nabla}f(1, 1, 1) = (0, 3, 4)$.

 C. Équation du plan tangent

Avec $x_0=1, y_0=1, z_0=1$ :
\begin{flalign*}
0(x - 1) + 3(y - 1) + 4(z - 1) &= 0\\
3y - 3 + 4z - 4 &= 0\\
3y + 4z &= 7
\end{flalign*}
\subsection{Déterminer le potentiel d'un champ de vecteurs conservatif}
\subsubsection{Champ de Vecteurs 2D : $\vec{F}(x,y)=(y-3x^2)\vec{e}_{x}+(x+\sin(y))\vec{e}_{y}$}
Un champ de vecteurs $\vec{F}(x,y) = P(x,y)\vec{e}_{x} + Q(x,y)\vec{e}_{y}$ défini sur le domaine simplement connexe $\mathbb{R}^2$ est conservatif (ou est un champ de gradient) si la condition de Schwarz est vérifiée :
$$
\frac{\partial Q}{\partial x} = \frac{\partial P}{\partial y}
$$

Ici, $P(x,y) = y-3x^2$ et $Q(x,y) = x+\sin(y)$, et

    $$
    \frac{\partial P}{\partial y} = \frac{\partial}{\partial y}(y-3x^2) = 1 = \frac{\partial Q}{\partial x} = \frac{\partial}{\partial x}(x+\sin(y))
    $$

Puisque $\frac{\partial Q}{\partial x} = \frac{\partial P}{\partial y} = 1$, le champ $\vec{F}$ est bien conservatif sur $\mathbb{R}^2$. On peut déterminer son potentiel :

Nous cherchons une fonction scalaire $f(x,y)$ telle que :
$$
\frac{\partial f}{\partial x} = P(x,y) = y-3x^2 \quad \mathbf{(1)}
$$
$$
\frac{\partial f}{\partial y} = Q(x,y) = x+\sin(y) \quad \mathbf{(2)}
$$

1.  Intégration de (1) par rapport à $x$ :
    $$
    f(x,y) = \int (y-3x^2) dx = yx - x^3 + C_1(y)
    $$
    où $C_1(y)$ est une fonction arbitraire dépendant uniquement de $y$.

2.  Dérivation par rapport à $y$ et identification avec (2) :
    Nous dérivons le résultat obtenu par rapport à $y$ :
    $$
    \frac{\partial f}{\partial y} = \frac{\partial}{\partial y}(yx - x^3 + C_1(y)) = x + C_1'(y)
    $$
    En comparant avec l'équation $\mathbf{(2)}$ ($\frac{\partial f}{\partial y} = x+\sin(y)$) :
    $$
    x + C_1'(y) = x + \sin(y)
    $$
    $$
    C_1'(y) = \sin(y)
    $$

3.  Intégration de $C_1'(y)$ pour trouver $C_1(y)$ :
    $$
    C_1(y) = \int \sin(y) dy = -\cos(y) + C
    $$
    où $C$ est une constante réelle.

4.  Conclusion :
    En substituant $C_1(y)$ dans l'expression de $f(x,y)$ on trouve l'expression finale de f :
    $$
    f(x,y) = yx - x^3 - \cos(y) + C
    $$

\subsubsection{Champ de Vecteurs 3D : $\vec{F}(x,y,z)=y^2\vec{e}_{x}+(2xy+z)\vec{e}_{y}+(y+3)\vec{e}_{z}$}

Un champ de vecteurs $\vec{F}(x,y,z) = P\vec{e}_{x} + Q\vec{e}_{y} + R\vec{e}_{z}$ est conservatif sur $\mathbb{R}^3$ si son rotationnel est nul : $\vec{\nabla} \wedge \vec{F} = \vec{0}$.
Les trois conditions à vérifier sont :
$$
\frac{\partial R}{\partial y} = \frac{\partial Q}{\partial z}, \quad \frac{\partial P}{\partial z} = \frac{\partial R}{\partial x}, \quad \frac{\partial Q}{\partial x} = \frac{\partial P}{\partial y}
$$
Ici, $P = y^2$, $Q = 2xy+z$, $R = y+3$.

On a bien :

\begin{flalign*}
\frac{\partial R}{\partial y} &= \frac{\partial}{\partial y}(y+3) &= &1 = \frac{\partial}{\partial z}(2xy+z) &= \frac{\partial Q}{\partial z}\\
\frac{\partial P}{\partial z} &= \frac{\partial}{\partial z}(y^2) &= &0 = \frac{\partial}{\partial x}(y+3) &= \frac{\partial R}{\partial x}\\
\frac{\partial Q}{\partial x} &= \frac{\partial}{\partial x}(2xy+z) &= &2y =  \frac{\partial}{\partial y}(y^2) &= \frac{\partial P}{\partial y}
\end{flalign*}

Le champ $\vec{F}$ est bien conservatif sur $\mathbb{R}^3$.

On peut donc déterminer le potentiel $f$ tel que $\vec{F} = \vec{\nabla}f$

Nous cherchons $f(x,y,z)$ telle que :
\begin{flalign*}
\frac{\partial f}{\partial x} &= y^2 \quad \mathbf{(1)}\\
\frac{\partial f}{\partial y} &= 2xy+z \quad \mathbf{(2)}\\
\frac{\partial f}{\partial z} &= y+3 \quad \mathbf{(3)}
\end{flalign*}

1.  Intégration de (1) par rapport à $x$ :
    $$
    f(x,y,z) = \int y^2 dx = y^2x + C_1(y,z)
    $$

2.  Dérivation par rapport à $y$ et identification avec (2) :
    $$
    \frac{\partial f}{\partial y} = \frac{\partial}{\partial y}(y^2x + C_1(y,z)) = 2yx + \frac{\partial C_1}{\partial y}
    $$
    En comparant avec $\mathbf{(2)}$ ($\frac{\partial f}{\partial y} = 2xy+z$) :
    $$
    2yx + \frac{\partial C_1}{\partial y} = 2xy+z \quad \Rightarrow \quad \frac{\partial C_1}{\partial y} = z
    $$

3.  Intégration de $\frac{\partial C_1}{\partial y}$ par rapport à $y$ :
    $$
    C_1(y,z) = \int z dy = zy + C_2(z)
    $$
    où $C_2(z)$ est une fonction arbitraire dépendant uniquement de $z$.
    Donc, $f(x,y,z) = y^2x + zy + C_2(z)$.

4.  Dérivation par rapport à $z$ et identification avec (3) :
    $$
    \frac{\partial f}{\partial z} = \frac{\partial}{\partial z}(y^2x + zy + C_2(z)) = y + C_2'(z)
    $$
    En comparant avec $\mathbf{(3)}$ ($\frac{\partial f}{\partial z} = y+3$) :
    $$
    y + C_2'(z) = y + 3 \quad \Rightarrow \quad C_2'(z) = 3
    $$

5.  Intégration de $C_2'(z)$ pour trouver $C_2(z)$ :
    $$
    C_2(z) = \int 3 dz = 3z + C
    $$
    où $C$ est une constante réelle.

6.  Conclusion :
    En substituant $C_2(z)$ dans l'expression de $f(x,y,z)$, on trouve :
    $$
    f(x,y,z) = xy^2 + yz + 3z + C
    $$
\subsection{Utiliser la formule de Green}
Le champ est $\vec{F}(x,y) = P\vec{e}_x + Q\vec{e}_y$ avec $P(x,y) = -y^3$ et $Q(x,y) = 2$.
La courbe $C$ est le cercle $x^2 + y^2 = 4$ (rayon $R=2$). $\mathcal{R}$ est le disque $x^2 + y^2 \le 4$.
\subsubsection{Calcul direct}
On calcule d'abord circulation $\int_{\mathcal{C}}\vec{F}\cdot d\vec{r}$

Paramétrisation de $C$ :
$$
x = 2\cos t, \quad y = 2\sin t, \quad 0 \le t \le 2\pi
$$
$$
dx = -2\sin t \dd t, \quad dy = 2\cos t \dd t
$$
Intégrale :
\begin{flalign*}
\int_{\mathcal{C}}\vec{F}\cdot d\vec{r} &= \int_{0}^{2\pi} [P dx + Q dy]\\
&= \int_{0}^{2\pi} [(-y^3)(-2\sin t) + (2)(2\cos t)] \dd t\\
&= \int_{0}^{2\pi} [-(2\sin t)^3 (-2\sin t) + 4\cos t] \dd t\\
&= \int_{0}^{2\pi} [16\sin^4 t + 4\cos t] \dd t
\end{flalign*}
L'intégrale de $4\cos t$ sur $[0, 2\pi]$ est nulle. Pour $\sin^4 t$, on utilise l'identité $\int_{0}^{2\pi} \sin^4 t \dd t = \frac{3\pi}{4}$.
$$
\int_{\mathcal{C}}\vec{F}\cdot d\vec{r} = 16 \cdot \frac{3\pi}{4} + 0 = 12\pi
$$
\subsubsection{Utilisation de Green}
On calcule l'intégrale double $\iint_{\mathcal{R}}(\frac{\partial Q}{\partial x}-\frac{\partial P}{\partial y})dx dy$

Rotationnel 2D :
\begin{flalign*}
\frac{\partial Q}{\partial x} &= \frac{\partial}{\partial x}(2) = 0\\
\frac{\partial P}{\partial y} &= \frac{\partial}{\partial y}(-y^3) = -3y^2\\
\frac{\partial Q}{\partial x}-\frac{\partial P}{\partial y} &= 0 - (-3y^2) = 3y^2
\end{flalign*}
Intégrale double (en coordonnées polaires) : $x=\rho \cos \varphi$, $y=\rho \sin \varphi$, $\dd A = \rho \dd\rho \dd\varphi$. On a
\begin{flalign*}
\iint_{\mathcal{R}} 3y^2 \dd x \dd y &= \int_{0}^{2\pi} \int_{0}^{2} 3(\rho \sin \varphi)^2 \rho \dd\rho \dd\varphi\\
&= 3 \left(\int_{0}^{2} \rho^3 \dd \rho\right) \left(\int_{0}^{2\pi} \sin^2 \varphi \dd \varphi\right)\\
&= 3 \cdot \left[\frac{\rho^4}{4}\right]_{0}^{2} \cdot \left[\pi\right]\\
&=3 \cdot \left(\frac{16}{4}\right) \cdot \pi = 3 \cdot 4 \cdot \pi = 12\pi
\end{flalign*}
On trouve bien le m\^eme résultat avec les 2 méthodes.
 \end{document}
